\chapter{API Loading and Error States}

Every component that fetches data moves through the same states, whether
it's a task list, a profile, or a dashboard widget. Handling them
consistently is what makes an app feel solid on a slow network.

\section{The State Chain}

A data-fetching component moves through \texttt{idle}, \texttt{loading},
then \texttt{success} or \texttt{error}. It should never leave the user
staring at a blank screen or a frozen spinner.

\begin{diagrambox}[API Component State Chain]
\begin{verbatim}
 idle -> loading (spinner/skeleton) -> success or error
                                          |         |
                              empty? -----+         v
                              yes: empty state   error msg + Retry
                              no: render list     -> Retry -> loading
\end{verbatim}
\end{diagrambox}

\begin{notebox}[NOTE]
There are really four outcomes: loading, error, empty (success, no data),
and populated. Most bugs come from forgetting the empty state.
\end{notebox}

\section{Implementing the Pattern}

Three state variables --- data, loading, error --- cover almost every
fetching component.

\begin{lstlisting}[language=JavaScript]
// components/TaskList.jsx
import { useEffect, useState } from "react";
import api from "../api/axios";

function TaskList() {
  const [tasks, setTasks] = useState([]);
  const [loading, setLoading] = useState(true);
  const [error, setError] = useState(null);

  const fetchTasks = async () => {
    setLoading(true);
    setError(null);
    try {
      const res = await api.get("/tasks");
      setTasks(res.data.data);
    } catch (err) {
      setError(err.response?.data?.message || "Failed to load tasks");
    } finally {
      setLoading(false);
    }
  };

  useEffect(() => { fetchTasks(); }, []);

  if (loading) return <Spinner />;
  if (error) {
    return (
      <div className="text-center py-8">
        <p className="text-red-600 mb-3">{error}</p>
        <button onClick={fetchTasks} className="px-4 py-2 bg-blue-600 text-white rounded">
          Retry
        </button>
      </div>
    );
  }
  if (!tasks.length) {
    return <p className="text-gray-500 text-center py-8">No tasks yet. Create your first one!</p>;
  }

  return (
    <ul className="space-y-2">
      {tasks.map((task) => <li key={task._id}>{task.title}</li>)}
    </ul>
  );
}

export default TaskList;
\end{lstlisting}

\begin{warnbox}[WARNING]
Always reset \texttt{error} to \texttt{null} at the start of a new fetch,
or a successful retry can leave a stale error message on screen.
\end{warnbox}

\section{Loading Spinners}

A spinner suits fast, small requests where showing layout structure would
be overkill.

\begin{lstlisting}[language=JavaScript]
// components/Spinner.jsx
function Spinner() {
  return (
    <div className="flex justify-center items-center py-8">
      <div className="h-8 w-8 border-4 border-blue-600 border-t-transparent
                       rounded-full animate-spin" />
    </div>
  );
}

export default Spinner;
\end{lstlisting}

\section{Skeleton Screens}

For content-heavy views, a skeleton that mimics the eventual layout
reduces perceived loading time better than a spinner.

\begin{lstlisting}[language=JavaScript]
// components/TaskCardSkeleton.jsx
function TaskCardSkeleton() {
  return (
    <div className="border rounded-lg p-4 space-y-3 animate-pulse">
      <div className="h-4 bg-gray-200 rounded w-3/4" />
      <div className="h-3 bg-gray-200 rounded w-1/2" />
      <div className="h-3 bg-gray-200 rounded w-full" />
    </div>
  );
}
// {loading ? Array.from({length:4}).map((_,i)=><TaskCardSkeleton key={i}/>) : tasks.map(t=><TaskCard key={t._id} task={t}/>)}
\end{lstlisting}

\begin{tipbox}[TIP]
Use spinners for quick actions, skeletons for initial page/list loads --
skeletons avoid the blank-then-full flash.
\end{tipbox}

\section{Choosing the Right State to Show}

\begin{center}
\begin{tabularx}{\textwidth}{lX}
\toprule
\textbf{Situation} & \textbf{What to Render} \\
\midrule
In flight & Spinner (actions) or skeleton (page loads) \\
Failed & Error message + Retry, never blank \\
Empty success & Friendly empty state with a call to action \\
Populated & The actual list/content \\
\bottomrule
\end{tabularx}
\end{center}

\whatwhyhow
{Every API-driven component tracks \texttt{loading}, \texttt{error}, and \texttt{data} state and renders a different UI for each.}
{Without explicit states, slow networks show blank screens and failures show nothing or crash.}
{Wrap the fetch in try/catch/finally, set the three variables, and branch render order: loading, error+Retry, empty, then populated.}
